% III. Програмна реализация.
% Приложението с изходен текст е свито до един указател към хранилището: цели
% файлове, вмъкнати в PDF, не се четат от никого.
\section{Програмна реализация}
\label{sec:implementation}

Реализацията е на C++20 и се изгражда с CMake. Разделението следва слоевете от
Раздел~\ref{sec:design}: договорите са в \code{include/lexicon/}, реализацията в
\code{src/}, данните в \code{data/}, тестовете в \code{tests/}. Всеки слой е отделна
единица за превод с един заглавен файл, тоест зависимостите са видими в списъка с
включвания и не могат да се образуват мълчаливо. Външни зависимости няма, и това е решение,
а не пропуск: изискването е чужд човек да може да изгради проекта на чиста машина. Затова
изпълнителят на тестовете е в самото хранилище, а базата от знания е текстов файл, четен
при стартиране.

Три решения заслужават коментар. \textbf{Първо}, един и същ възел от гората може да бъде
построен от повече от една продукция, и тогава завършването предлага една и съща връзка два
пъти; без проверка за повторение броят на изводите нараства, без изречението да е станало
по-нееднозначно. \textbf{Второ}, отхвърлянето на разчитане е нормална работа, не програмна
грешка: когато моделът няма рамка за предлога, причината се предава с изключение от отделен
тип, което конвейерът улавя и записва дословно, вместо да прекъсва анализа. \textbf{Трето},
нито един от двата генератора не вижда въпроса, дървото или лексикона; този на SQL цитира
променлива от обхващащо ниво, което прави вложеното \code{NOT EXISTS} корелирано.

Тестовете са 30 отделни случая в пет файла, всеки регистриран в CTest поотделно, така че
провалът назовава случая, а не файла. Че логическата форма е абстракция, а не преоблечен
шаблон, също е тест: активният и страдателният залог пораждат побайтово еднакви заявки и на
двата целеви езика. Пълният изходен текст, данните и записаните измервания са в публичното
хранилище \url{https://github.com/Alex-Tsvetanov/lexicon}; обемът, измерен с \code{wc -l},
е в Табл.~\ref{tab:repo}.

\begin{table}[H]
\caption{Състав на хранилището}
\label{tab:repo}
\centering
\begin{tabular}{@{}L{3.2cm}L{7.2cm}r@{}}
\toprule
\textbf{Път} & \textbf{Съдържание} & \textbf{Редове} \\
\midrule
\code{include/}, \code{src/} & по един договор и една реализация на слой & 3345 \\
\code{data/} & лексикони, граматики, схема, знания, въпроси & 473 \\
\code{apps/}, \code{tests/} & демонстрация, измерване, 30-те тестови случая & 1045 \\
\code{tools/}, \code{measurements/} & проверката на SQL и записаните измервания & 236 \\
\bottomrule
\end{tabular}
\end{table}
